\section{第2章 导数与微分}
\subsection{导数的基本概念}
\begin{exercise}{151}{定义求导与线性主部}
用导数定义求 $f(x)=1+2x+3x^2$ 的 $f'(1)$，并把 $f(1+h)-f(1)$ 写成线性主部与余项之和，验证余项为 $o(h)$。
\end{exercise}
\begin{exercise}{152}{零因子处的导数}
设 $f(x)=x(2x+1)(3x+2)(4x+3)$。不展开多项式，求 $f'(0)$；再求 $f'(-1/2)$，说明在单零点处如何利用其余因子的连续性。
\end{exercise}
\begin{exercise}{153}{定义求导与零点因子}
设 $f(x)=x(x-2)\sqrt{\frac{3x^2+1}{2x^2+1}}$。用导数定义求 $f'(0)$、$f'(2)$，并分别写出曲线在这两点的切线方程。
\end{exercise}
\begin{exercise}{154}{补定义函数的导数}
设 $f(x)=\begin{cases}\sin^2x/(2x),&x\ne0,\\0,&x=0.\end{cases}$ 用定义求 $f'(0)$；若把 $f(0)$ 改为参数 $c$，确定在 $0$ 处连续、可导各需要什么条件。
\end{exercise}
\begin{exercise}{155}{连续因子与乘积可导}
设 $f(x)=(x^2-1)\varphi(x-1)$，$\varphi$ 在 $0$ 处连续。证明 $f$ 在 $1$ 处可导并求 $f'(1)$；构造 $\varphi$ 在 $0$ 处不可导而该结论仍成立的例子。
\end{exercise}
\begin{exercise}{156}{平移函数关系的导数}
设 $f$ 在 $0$ 处可导，$f'(0)=1/2$，且对所有实数 $x$ 有 $f(2+x)=2f(x)$。用定义求 $f'(2)$；进一步证明 $f$ 在所有偶整数点可导，并求 $f'(2n)$（$n\in\mathbb Z$）。
\end{exercise}
\begin{exercise}{157}{差商极限与点值}
设 $f$ 在 $0$ 处连续，且 $\lim_{x\to0}\frac{f(x)+3}{x}=2$。判断 $f$ 在 $0$ 处是否可导并求导数；若去掉连续性，说明能否保留可导结论，并给出反例。
\end{exercise}
\begin{exercise}{158}{缩放增量的差商}
设 $f$ 在 $x_0$ 处可导。求 $\lim_{h\to0}\frac{f(x_0+5h)-f(x_0)}h$ 与 $\lim_{h\to0}\frac{f(x_0)-f(x_0-3h)}h$；用同一方法写出一般系数 $a,b$ 下 $\lim_{h\to0}\frac{f(x_0+ah)-f(x_0+bh)}h$。
\end{exercise}
\begin{exercise}{159}{双移位差商与逆命题}
设 $f'(x_0)$ 存在，求 $\lim_{h\to0}\frac{f(x_0+4h)-f(x_0+2h)}h$。若不要求 $f$ 在 $x_0$ 连续，构造一个在该点不可导但上述极限仍存在的函数，说明差分式中没有出现点值的影响。
\end{exercise}
\begin{exercise}{160}{二次增量下的单侧信息}
设 $f$ 在 $3$ 处可导，且 $\lim_{x\to0}\frac{f(3)-f(3-3x^2)}{\ln(1+x^2)}=3$。求 $f'(3)$；说明若仅假设 $f$ 在 $3$ 处连续，给定极限能否保证双侧导数存在。
\end{exercise}
\begin{exercise}{161}{可导函数的渐近形式}
设 $f(0)=0$、$f'(0)=1$。求 $\lim_{x\to0}\frac{f(3x)-f(-5x)}x$ 与 $\lim_{x\to0}\frac{f(2x^2)-xf(x)}{x^2}$；分别写明使用 $f(t)=t+o(t)$ 时余项的变化。
\end{exercise}
\begin{exercise}{162}{可导充分条件的辨析}
设 $f(0)=0$，分别判断下列极限存在是否足以保证 $f'(0)$ 存在，证明或举反例：$\lim_{h\to0}\frac{f(3h)-f(h)}{2h}$，$\lim_{h\to0}\frac{f(h)-f(-h)}{2h}$，$\lim_{h\to0}\frac{f(h^3)}{h^3}$，$\lim_{h\to0}\frac{f(h^2)}{h^2}$。指出奇次与偶次换元提供的自变量方向有何不同。
\end{exercise}
\begin{exercise}{163}{二次内增量与已知导数}
设 $f$ 在 $1$ 处可导，$f(1)=0$、$f'(1)=2$。求 $\lim_{x\to0}\frac{\ln(1+\sin^2(3x))-f(\cos x)}{x^2}$，写出 $f(\cos x)$ 的首个非零渐近项；说明为何只知道 $f'(1)$ 已足以计算该极限。
\end{exercise}
\begin{exercise}{164}{内点导数与端点约束}
分别判断 $|x|$、$\sqrt x$、$x^2$、$x\ln x$ 在 $0$ 处是否具有通常的双侧导数。对原式在 $0$ 未定义的函数，先判断能否连续补定义，再讨论其右差商是否有有限极限。
\end{exercise}
\begin{exercise}{165}{左右导数的拼接}
设 $f(x)=\begin{cases}\cos x,&x>0,\\x^2+1,&x\le0.\end{cases}$ 求 $f'_-(0)$、$f'_+(0)$ 并判断可导性；再比较两侧二阶导数，判断 $f''(0)$ 是否存在。
\end{exercise}
\begin{exercise}{166}{指数分段函数的单侧导数}
设 $f(x)=\begin{cases}x(1+e^{1/x}),&x\ne0,\\0,&x=0.\end{cases}$ 分别讨论其在 $0$ 处的左连续性、右连续性与左右导数，区分有限单侧导数与无穷差商；说明一侧导数存在能否推出另一侧连续。
\end{exercise}
\begin{exercise}{167}{绝对值造成的尖点}
讨论 $f(x)=x^2|x-1|$ 在 $1$ 处的左右导数与可导性；若将因子 $x^2$ 换为 $(x-1)^2$，重新判断并解释零点因子的作用。
\end{exercise}
\begin{exercise}{168}{零点消除绝对值尖点}
讨论 $f(x)=(x^3-1)|x-1|$ 在 $1$ 处的可导性并求导数；进一步判断 $f'$ 在 $1$ 处是否连续、$f''(1)$ 是否存在。
\end{exercise}
\begin{exercise}{169}{可导的充要条件}
设 $\varphi$ 在 $1$ 处连续，$f(x)=|x-1|\varphi(x)$。证明 $f$ 在 $1$ 处可导当且仅当 $\varphi(1)=0$，并求此时的导数；说明证明中连续性使用在何处。
\end{exercise}
\begin{exercise}{170}{绝对值复合的全部不可导点}
求 $f(x)=(x^2+x-2)|x^2-4x|$ 的全部不可导点及其左右导数；说明绝对值内部多项式的零点是否一定造成不可导。
\end{exercise}
\begin{exercise}{171}{切线与法线}
求曲线 $y=2x+\cos x$ 在 $(0,1)$ 处的切线、法线方程；再确定该曲线是否存在水平切线，并给出依据。
\end{exercise}
\begin{exercise}{172}{切线垂直的可行性}
在相同横坐标 $x_0$ 处，曲线 $y=x^2$ 与 $y=x^3$ 的切线互相垂直。求所有 $x_0$，并检查用斜率乘积为 $-1$ 的判据时是否遗漏竖直切线情形。
\end{exercise}
\begin{exercise}{173}{指定方向的切线}
直线 $L$ 与 $2x-6y+1=0$ 垂直，且与曲线 $y=x^3+3x^2-5$ 相切。求全部可能的切点及 $L$ 的方程；说明斜率方程的重根是否表示存在两条不同切线。
\end{exercise}
\begin{exercise}{174}{切线截距与面积不变量}
证明双曲线 $xy=1$ 上任一点处的切线与坐标轴围成的三角形面积恒为 $2$；写明两截距的符号，并说明点位于第三象限时面积结论为何仍成立。
\end{exercise}
\begin{exercise}{175}{公共切线的两个条件}
求 $a,b$，使 $y=1+\ln x$ 与 $y=x^2+ax+b$ 在 $x=1$ 处具有同一条切线；分别列出点值条件和斜率条件，并求公共切线方程。
\end{exercise}
\begin{exercise}{176}{偶函数导数与倒数差商}
设 $f$ 为可导偶函数，且 $\lim_{h\to0}\frac{h}{f(3+h)-f(3)}=-2$。求曲线在 $(-3,f(-3))$ 处的切线斜率；从奇偶性定义证明所使用的导数关系。
\end{exercise}
\begin{exercise}{177}{可导蕴含连续}
设 $f$ 在 $2$ 处可导，且 $\lim_{x\to2}f(x)=5$。求 $f(2)$；证明所使用的“可导蕴含连续”结论，并构造函数说明连续不足以保证可导。
\end{exercise}
\begin{exercise}{178}{振荡函数的连续与可导}
设 $f(x)=\begin{cases}x\cos(1/x^2),&x\ne0,\\0,&x=0.\end{cases}$ 讨论 $0$ 处的连续性与可导性；若将 $x$ 改为 $x^2$，重新讨论，并用点列解释差商无极限的情形。
\end{exercise}
\begin{exercise}{179}{连续可导的分段拼接}
设 $f(x)=\begin{cases}xe^x,&x\ge0,\\x^3+x,&x<0.\end{cases}$ 判断其在 $0$ 处的连续性、可导性和导函数连续性；进一步判断 $f''(0)$ 是否存在。
\end{exercise}
\begin{exercise}{180}{先连续再比较导数}
设 $f(x)=\begin{cases}x^3+\ln x,&x\ge1,\\x^3\ln x,&0<x<1.\end{cases}$ 判断 $1$ 处的连续性与可导性；若只允许给右支加一个常数 $c$，确定能使函数连续的 $c$，并判断此时是否可导。
\end{exercise}
\begin{exercise}{181}{可导拼接反求参数}
求 $a,b$，使 $f(x)=\begin{cases}x\sin\pi x,&x>1,\\ax^2+b,&x\le1\end{cases}$ 在 $1$ 处可导；再判断此时 $f'$ 在 $1$ 处是否连续及 $f''(1)$ 是否存在。
\end{exercise}
\begin{exercise}{182}{有界因子能提供的局部信息}
设 $\varphi$ 在 $0$ 的某邻域内有界，$f(x)=\begin{cases}\ln(1+2x^2)/x^2,&x>0,\\x^2\varphi(x),&x\le0.\end{cases}$ 求 $0$ 处左右极限并判断连续性、可导性；说明给 $f(0)$ 重新赋值能否修复该处的间断。
\end{exercise}
\begin{exercise}{183}{对称差商的条件辨析}
设 $f$ 在 $0$ 处连续，且 $\lim_{h\to0}\frac{f(h)-f(-h)}{2h}=L$。举例说明 $f'(0)$ 未必存在；若再有 $\lim_{h\to0}\frac{f(h)+f(-h)-2f(0)}h=0$，证明 $f'(0)=L$。
\end{exercise}
\begin{exercise}{184}{绝对值复合的可导判据}
设 $g$ 在 $x_0$ 处可导且 $g(x_0)=0$。证明 $|g|$ 在 $x_0$ 处可导当且仅当 $g'(x_0)=0$；再求 $|x^3-3x+2|$ 的全部不可导点，并解释其中重根与单根的区别。
\end{exercise}
\subsection{函数的求导法则}
\begin{exercise}{185}{基本导数与定义域}
设 $f(x)=3x^2+2e^x-5\ln x+\sin1$。求 $f'(x)$、$f'(1)$ 及 $x=1$ 处的切线方程，说明常数项在点值和斜率中的不同作用。
\end{exercise}
\begin{exercise}{186}{乘积法则与定点导数}
设 $f(x)=x\cos x+3x^2$。求 $f'(x)$、$f'(\pi/2)$；判断表达式 $(x\cos x)'=\cos x\,(x)'$ 是否成立，并写明遗漏项。
\end{exercise}
\begin{exercise}{187}{指数三角乘积的求导}
设 $y=e^x(\cos x-2\sin x)$。求 $y'$，并求 $x\in[0,2\pi)$ 中全部水平切线所对应的横坐标。
\end{exercise}
\begin{exercise}{188}{对数因子的定义范围}
设 $y=(1+x^2)\arctan x-x\ln x$。求 $y'$；将函数在 $0$ 处连续补定义后，判断其右导数是否为有限值。
\end{exercise}
\begin{exercise}{189}{商式化简后求导}
设 $y=\frac{\sqrt x+xe^x-3x^3\sin x}{x^2}$。写出定义域并求 $y'$；比较直接使用商法则与先拆成三项时的表达式，验证结果一致。
\end{exercise}
\begin{exercise}{190}{商法则与分母零点}
设 $y=\frac{\sin x}{x-\cos x}$。求 $y'$ 并写明适用条件；证明方程 $x-\cos x=0$ 只有一个实根，判断该点能否连续补定义。
\end{exercise}
\begin{exercise}{191}{对数有理式的求导}
设 $y=\frac{2\ln x+x^3}{3\ln x+x^2}$。求 $y'$、$y'(1)$，并说明导数表达式使用时对 $x$ 及分母的限制；证明正半轴上恰有一个分母零点。
\end{exercise}
\begin{exercise}{192}{指数有理式的乘商求导}
设 $y=\frac{e^x(x^2+1)}{x+1}$。求 $y'$ 并化简，写出定义域；进一步求 $\lim_{x\to-1}(x+1)^2y'$，说明奇点处不能直接代入求导公式。
\end{exercise}
\begin{exercise}{193}{换元确定函数及导数}
已知对一切 $x\ne0$，$f(x-1/x)=\frac{x^4+1}{x^4+x^2+1}$。求 $f(t)$、$f'(t)$，说明换元后 $t$ 的范围；验证不同原像给出的函数值一致。
\end{exercise}
\begin{exercise}{194}{分段函数的导函数}
设 $f(x)=\begin{cases}x/(2x-1),&x<0,\\e^x\sin x,&x\ge0.\end{cases}$ 求所有可导点处的 $f'(x)$，单独判断 $0$ 处是否可导，并写出该处左右导数。
\end{exercise}
\begin{exercise}{195}{反函数导数的取点}
设可导函数 $f$ 的反函数为 $\varphi$，已知 $f(1)=2$、$f(2)=3$、$f'(1)=4$、$f'(2)=5$。求 $\varphi'(2)$、$\varphi'(3)$，说明为何不能将 $\varphi'(2)$ 直接写为 $1/f'(2)$。
\end{exercise}
\begin{exercise}{196}{反函数存在性与导数}
设 $f(x)=3x^3+2x+1$。证明其在 $\mathbb R$ 上有反函数 $\varphi$，求 $\varphi'(1)$；进一步求 $\varphi''(1)$，并说明所用分母为何非零。
\end{exercise}
\begin{exercise}{197}{多层复合的链式法则}
设 $y=(1-x+2x^3)^4$。求 $y'(0)$、$y'(1)$，并写出导函数；比较内函数在两点处的值及其导数各如何进入结果。
\end{exercise}
\begin{exercise}{198}{对数三角指数复合}
设 $y=\ln(\cos(e^{-x}))$。求 $y'$，写出函数的完整定义域；确定包含 $x=0$ 的最大定义区间，并说明该区间端点处表达式的变化。
\end{exercise}
\begin{exercise}{199}{根式反三角函数的求导范围}
设 $f(x)=\arcsin\sqrt{\frac{1-x}{1+x}}$。求定义域、内部各点的 $f'(x)$ 及 $f'(1/2)$；判断定义域两端点的单侧差商是否有有限极限。
\end{exercise}
\begin{exercise}{200}{共轭式化简与端点导数}
设 $y=\frac{\sqrt{x+1}-\sqrt{x-1}}{\sqrt{x+1}+\sqrt{x-1}}$。求定义域与内部导数，给出一种先化简再求导的方法；判断 $x=1$ 处右差商是否有有限极限。
\end{exercise}
\begin{exercise}{201}{多种指数结构的求导}
设 $y=3^x+x^3+x^{2x}+3^{x^2}$（$x>0$）。求 $y'$、$y'(1)$；分别指出常底指数、常指数幂、幂指函数及复合指数项所用的求导规则。
\end{exercise}
\begin{exercise}{202}{三角函数乘积与复合}
设 $y=\cos3x\,\sin^3x$。求 $y'$ 与 $y'(\pi/6)$；将 $y'$ 在 $\sin x=0$ 和 $\cos3x=0$ 的点分别代入，说明不能通过直接除去这些因子求全部驻点。
\end{exercise}
\begin{exercise}{203}{反三角函数与根式导数}
设 $y=x\sqrt{1-x^2}-\arccos x$。求内部导数并尽量化简，写出适用区间；用差商判断两个端点的单侧导数是否有限。
\end{exercise}
\begin{exercise}{204}{根式与对数的协同化简}
设 $y=x^2+x\sqrt{x^2+1}+\ln(x+\sqrt{x^2+1})$。证明对数对所有实数 $x$ 有定义，求 $y'$ 并化简；据此判断该函数在 $\mathbb R$ 上的单调性。
\end{exercise}
\begin{exercise}{205}{对数相位的复合导数}
设 $y=x[\cos(\ln x)+\sin(\ln x)]$。求 $y'$；将 $y$ 在 $0$ 处连续延拓后，判断右导数是否存在，并给出两列差商证明。
\end{exercise}
\begin{exercise}{206}{反余切指数复合求导}
设 $y=-x+\frac12\ln(1+e^{2x})-e^{-x}\operatorname{arccot}e^x$，其中 $\operatorname{arccot}\in(0,\pi)$。求 $y'$ 并化简，再判断该函数在 $\mathbb R$ 上的单调性。
\end{exercise}
\begin{exercise}{207}{先取参数极限再求导}
对每个固定实数 $x$，设 $f(x)=\lim_{t\to+\infty}x\left(\frac{t-x}{t+2x}\right)^t$，实数幂在充分大的 $t$ 上考虑。求 $f(x)$ 及 $f'(x)$；先计算定义函数的极限，不得未经论证交换求导与取极限。
\end{exercise}
\begin{exercise}{208}{抽象函数的链式法则}
设 $f$ 可导。分别求 $y=f(3^x+x^2)$ 与 $y=f(x\ln x)$ 的导数，写明第二式的定义范围；若 $f'(1)=2$，再求第一式在 $x=0$ 处的导数。
\end{exercise}
\begin{exercise}{209}{复合求导的取点对应}
已知 $f$ 可导且 $f'(1)=-2$，$g(x)=f(2x^3-3x)$。求 $g'(-1)$；再求满足 $2x^3-3x=1$ 的全部实数 $x$，指出仅用已知的 $f'(1)$ 还可确定哪些点的 $g'$。
\end{exercise}
\begin{exercise}{210}{反函数与复合函数求导}
设 $f(x)=x+e^x$，其反函数为 $\varphi$。证明反函数在 $\mathbb R$ 上存在且可导，求 $\varphi'(1)$；再设 $G(x)=\varphi(e^x)$，求 $G'(0)$ 与 $G''(0)$。
\end{exercise}
\begin{exercise}{211}{函数关系的求导约束}
设 $f:\mathbb R\to\mathbb R$ 可导，满足 $f(x+y)=f(x)+f(y)+2xy$。先求 $f(0)$，再用固定 $x$、令增量趋零的方法证明 $f'(x)=f'(0)+2x$；若 $f'(0)=3$，确定 $f(x)$ 并验证。
\end{exercise}
\subsection{高阶导数}
\begin{exercise}{212}{二阶导数与定义域}
分别求 $y=\frac{\sin x-\cos x}{\sin x+\cos x}$ 与 $y=2x^3\arctan x+\ln(1+x^2)$ 的二阶导数；写明前者的定义范围，并求后者的 $y''(0)$。
\end{exercise}
\begin{exercise}{213}{指数三角函数的二阶导数}
设 $f(x)=e^{3x}\sin2x$。求 $f''(\pi/2)$；进一步找出常数 $a,b$，使 $f''+af'+bf\equiv0$，并直接验证。
\end{exercise}
\begin{exercise}{214}{对数根式的二阶求导}
设 $y=\ln\sqrt{\frac{1+e^{2x}}{e^{2x}}}$。求 $y'(0)$、$y''(0)$，再写出 $y$ 在 $0$ 处到二阶的展开式，并说明余项形式。
\end{exercise}
\begin{exercise}{215}{反正弦函数的二阶导数}
设 $y=\arcsin(x/3)$。求 $y''(\sqrt5)$ 及 $y''$ 的定义区间；根据 $y''$ 的符号判断其在负、正半区间上的凹凸性。
\end{exercise}
\begin{exercise}{216}{对数乘积的高阶导数}
设 $y=x\ln x$（$x>0$）。求 $y^{(10)}$，再给出 $y^{(n)}$（$n\ge2$）的通式并用数学归纳法证明，明确符号与阶乘的下标。
\end{exercise}
\begin{exercise}{217}{三角函数导数的周期性}
设 $y=(\sin(x/2)+\cos(x/2))^2$。先化简再求 $y^{(10)}$；给出 $y^{(n)}$（$n\ge1$）的统一公式，并核对前两阶。
\end{exercise}
\begin{exercise}{218}{莱布尼茨公式的简化}
设 $f(x)=xe^{-x}$。求 $f^{(n)}(x)$（$n\ge1$）并证明；再求 $f^{(n)}(n)$，说明乘积中仅有两项对高阶导数有贡献的原因。
\end{exercise}
\begin{exercise}{219}{多项式次数与最高阶导数}
设 $y=(x^2+1)^5(x^5+1)^2$。不完全展开，求 $y^{(20)}$、$y^{(21)}$；再求 $y^{(19)}(0)$，说明只需检查哪些幂次。
\end{exercise}
\begin{exercise}{220}{降幂后的任意阶导数}
设 $f(x)=\sin^4x+\cos^4x$。化简并求 $f^{(n)}(x)$（$n>1$）；再求满足 $f^{(n)}(0)=0$ 的全部正整数 $n$。
\end{exercise}
\begin{exercise}{221}{复合对数的高阶导数}
设 $f(x)=\ln(2-3x)$。求 $f^{(n)}(x)$（$n\ge1$）并用归纳法证明，写出适用区间；再求 $f^{(n)}(0)$。
\end{exercise}
\begin{exercise}{222}{导数函数关系的继续求导}
设 $f$ 在 $1$ 的某邻域内可导，满足 $f'=f^2$ 且 $f(1)=2$。证明条件保证 $f''$ 存在，求 $f''(1)$；再求 $f'''(1)$，说明每次求导的合法性。
\end{exercise}
\begin{exercise}{223}{高阶零点与莱布尼茨公式}
设 $g$ 在 $a$ 附近有 $n+1$ 阶导数，$F(x)=(x-a)^ng(x)$，其中 $n\ge1$。求 $F^{(k)}(a)$（$0\le k\le n+1$），并用莱布尼茨公式证明；据此求 $(x-1)^5\ln x$ 在 $1$ 处的六阶导数。
\end{exercise}
\begin{exercise}{224}{高阶导数的递推公式}
设 $f(x)=e^{ax}\sin bx$。证明存在数列 $A_n,B_n$，使 $f^{(n)}(x)=e^{ax}(A_n\sin bx+B_n\cos bx)$，写出初值与递推式；取 $a=b=1$，求 $f^{(8)}(0)$，并用相位表示式核对。
\end{exercise}
\subsection{隐函数及由参数方程所确定的函数的导数}
\begin{exercise}{225}{隐函数求导与正则点}
分别对 $\cos(xy)+\ln y-x=1$、$x^2+y^2=6xy$ 求 $\dd y/\dd x$，写出公式分母非零的条件；对第二条曲线单独讨论原点的直线分支与切线。
\end{exercise}
\begin{exercise}{226}{隐函数分支与定点导数}
方程 $xy-\sin(\pi y^2)=0$ 在 $(0,1)$ 附近确定分支 $y=y(x)$。求 $y'(0)$；若改取过 $(0,\sqrt2)$ 的分支，重新计算并说明只给 $x=0$ 是否足以唯一确定导数。
\end{exercise}
\begin{exercise}{227}{隐曲线的切线与法线}
求曲线 $\tan(x+y+\pi/4)=y$ 在 $(-1,1)$ 处的切线与法线方程；检查该点满足原方程，并说明求导中 $y$ 对 $x$ 的依赖如何进入。
\end{exercise}
\begin{exercise}{228}{同一横坐标的多条隐分支}
求曲线 $y^2+2e^{xy}=3y$ 在横坐标 $x=0$ 对应的全部点及各点的切线方程；检查每一点处隐函数求导式的分母是否非零。
\end{exercise}
\begin{exercise}{229}{抽象隐函数的定点导数}
设 $xf(y)=e^y\ln y$ 确定过 $(0,1)$ 的可导分支，$f$ 可导且 $f(1)=1$。求 $y'(0)$；若再知 $f$ 在 $1$ 附近连续可导且 $f'(1)=c$，求 $y''(0)$，写明所需的连续性与可导性来自何处。
\end{exercise}
\begin{exercise}{230}{隐显曲线公共切线}
曲线 $y=x^2+ax+b$ 与 $2y=xy^2-3$ 在 $(1,-1)$ 处有公共切线。求全部参数 $a,b$ 及公共切线方程；分别写明通过指定点和切线斜率相同给出的条件。
\end{exercise}
\begin{exercise}{231}{隐函数的二阶导数}
在 $x\ne0$ 的区域，方程 $\ln(x^2+y^2)=2\arctan(y/x)+\pi/2+\ln2$ 确定局部可导分支 $y(x)$。求 $y'$、$y''$，写出导数公式的分母条件，并说明原方程的定义范围。
\end{exercise}
\begin{exercise}{232}{隐函数二阶导数的定点计算}
设 $e^y+xy=e+1$ 确定过 $(1,1)$ 的局部分支。求 $y'(1)$、$y''(1)$；写出该点处的一阶线性近似，并检查一阶求导式的分母非零。
\end{exercise}
\begin{exercise}{233}{含导数关系的再次求导}
设曲线 $y(x)$ 二阶可导，经过 $(1,-1)$，满足 $x^2+y^2+(y')^3=1$。求 $y'(1)$、$y''(1)$；说明为何能在该点附近由再次求导的方程解出 $y''$。
\end{exercise}
\begin{exercise}{234}{对数求导与正性}
设 $y=x^{\ln x}$（$x>0$）。用对数求导法求 $y'$，确定全部驻点及单调区间；说明取自然对数的合法性。
\end{exercise}
\begin{exercise}{235}{变底变指数的求导}
在 $1+\sin x>0$ 的区间内，设 $y=(1+\sin x)^{x^2}$。求 $y'$ 与 $y'(0)$；说明底数等于零的点为何不能直接代入对数求导公式。
\end{exercise}
\begin{exercise}{236}{幂指函数与根式乘积}
设 $y=x^{\sin x}\sqrt{e^x-1}$（$x>0$）。求 $y'$，再求 $\lim_{x\to0^+}y/\sqrt x$；据此判断将 $y(0)$ 定义为 $0$ 后右导数是否有限。
\end{exercise}
\begin{exercise}{237}{变指数的商函数}
设 $y=(1+2x)^{1/x^2}e^{-x}$，其中 $x>-1/2$、$x\ne0$。求 $y'$；分别讨论 $x\to0^+$、$x\to0^-$ 时 $y$ 的变化，判断能否在 $0$ 处连续补定义。
\end{exercise}
\begin{exercise}{238}{多因子对数求导}
设 $y=\frac{(x+1)^2\sqrt{2x+2}}{(x-1)^3\sqrt{3x-3}}$。确定实定义域，用对数求导法求 $y'$ 并化简；根据 $y'/y$ 的符号判断单调性。
\end{exercise}
\begin{exercise}{239}{幂指函数的一阶导数}
求 $y=x^x$（$x>0$）的 $y'$，确定其最小值及取值点；再求 $\lim_{x\to0^+}x^x$，判断连续延拓后的右导数是否有限。
\end{exercise}
\begin{exercise}{240}{幂指函数的二阶导数}
设 $y=x^{2x}$（$x>0$）。求 $y''$，并判断 $y''$ 在定义域内的符号；给出判断中使用的不等式或平方分解。
\end{exercise}
\begin{exercise}{241}{隐式幂指方程的分支}
在 $x,y>0$ 且 $1+\ln y\ne0$ 的分支上，方程 $x^{1/y}=y^{1/x}$ 确定 $y=y(x)$。求 $y'$、$y''$；说明取对数后方程如何化简，并核对分支 $y=x$ 上的结果。
\end{exercise}
\begin{exercise}{242}{参数方程求导与例外参数}
分别求参数曲线 $x=2t/(1+t^2),\ y=2t^2/(1+t^2)$ 和 $x=\cos^3t,\ y=\sin^3t$ 的 $\dd y/\dd x$。列出 $\dd x/\dd t=0$ 的参数，并对第一条曲线在这些参数处判断切线方向。
\end{exercise}
\begin{exercise}{243}{参数导数的定点取值}
设 $x=t^2e^{2t}$、$y=(t^2+1)e^t$。求 $\left.\dd y/\dd x\right|_{t=2}$；列出一阶参数求导公式不能直接使用的全部 $t$，并说明原因。
\end{exercise}
\begin{exercise}{244}{参数曲线切线与双分支}
求 $x=\sqrt{t^2+3}$、$y=\ln(t+\sqrt{t^2+3})$ 在 $t=1$ 对应点的切线方程；再比较 $t=-1$ 对应点的坐标及切线斜率，说明相同 $x$ 是否对应相同 $y$。
\end{exercise}
\begin{exercise}{245}{由坐标确定参数}
求曲线 $x=e^t\sin2t$、$y=e^t\cos t$ 在 $(0,1)$ 处的切线和法线方程；先确定该点对应的全部实数参数，检查是否存在多条切线。
\end{exercise}
\begin{exercise}{246}{参数曲线二阶导数}
设 $x=\cos t$、$y=t\cos t-\sin t$。求 $\dd y/\dd x$ 与 $\dd^2y/\dd x^2$，写出公式适用的参数范围；说明求二阶导数时为何还须除以 $\dd x/\dd t$。
\end{exercise}
\begin{exercise}{247}{反三角参数的二阶导数}
设 $x=\arctan t$、$y=3t+t^3$。求 $\left.\dd^2y/\dd x^2\right|_{t=1}$；消去参数写出显函数，并用显函数求导核对。
\end{exercise}
\begin{exercise}{248}{参数导数与全局单值性}
设 $x=e^t+t$、$y=\sin t$。求 $t=0$ 时的一阶、二阶 $x$ 导数；证明参数式在整个实轴上确定单值函数 $y(x)$，并写出其 $x$ 定义域。
\end{exercise}
\begin{exercise}{249}{隐函数正则点与退化点}
考虑曲线 $x^2+xy+y^2=3$。求过 $(1,1)$ 的局部分支的一阶、二阶导数及切线；找出曲线上 $\dd y/\dd x$ 公式分母为零的点，并判断这些点的切线方向。
\end{exercise}
\begin{exercise}{250}{参数退化点的导数定义}
设 $x=t^3$、$y=t^4$。消去参数并讨论原点处的可导性与切线，再讨论二阶导数是否存在；解释为什么 $t=0$ 时参数导数商为 $0/0$ 并不能直接判定不可导。
\end{exercise}
\subsection{函数的微分}
\begin{exercise}{251}{微分与增量的误差}
设 $y=x^3+x$。在 $x=3$ 处，分别取 $\Delta x=0.1,0.01$，计算 $\dd y$、$\Delta y$ 及二者之差；写出一般增量 $h$ 时的误差表达式，并说明误差相对于 $h$ 的阶。
\end{exercise}
\begin{exercise}{252}{可微定义与余项}
已知 $f$ 在每个 $x>-1$ 处的增量可写为 $\Delta y=\frac{\Delta x}{\sqrt{1+x^3}}+\alpha$，其中 $\alpha/\Delta x\to0$（$\Delta x\to0$）。求 $f'(2)$ 及 $\dd y|_{x=2}$；说明线性系数唯一的原因。
\end{exercise}
\begin{exercise}{253}{微分和增量的等价关系}
设 $f$ 在 $x_0$ 处可导且 $f'(x_0)=1$。比较 $\dd y$、$\Delta y$ 与 $\Delta x$ 在 $\Delta x\to0$ 时的等价关系并证明；若 $f'(x_0)=0$，说明还能保留哪些高阶关系。
\end{exercise}
\begin{exercise}{254}{复合函数的微分}
求 $y=e^{\sin^2(1/x)}$ 与 $y=\arcsin(x/2)+\sqrt{4-x^2}$ 的微分，并写出公式成立的范围；对第二个函数判断能否把内部微分公式直接用于端点。
\end{exercise}
\begin{exercise}{255}{定点微分的系数}
分别求 $y=\sqrt2+\sin x$ 在 $x=0$ 处和 $y=\ln(1+2^{-2x})$ 在 $x=1$ 处的微分；将 $\dd x=0.01$ 代入，写出相应的函数增量线性近似。
\end{exercise}
\begin{exercise}{256}{隐函数微分的唯一性}
设 $2y+\sin(x+y)=0$ 确定 $y=y(x)$。求 $\dd y$，证明导数公式的分母恒不为零；在过原点的分支上，求 $x=0$ 处的线性近似。
\end{exercise}
\begin{exercise}{257}{隐函数的定点微分}
设 $e^y+xy-e=0$ 在 $x=0$ 附近确定分支 $y(x)$。求该点对应的 $y$ 值及 $\dd y|_{x=0}$；据此估计 $x$ 增加 $0.01$ 时 $y$ 的增量，说明这是线性近似。
\end{exercise}
\begin{exercise}{258}{可去点处的微分}
设 $f(x)=\begin{cases}(\cos x-1)/x,&x\ne0,\\0,&x=0.\end{cases}$ 用定义求 $\dd f|_{x=0}$；说明为何不能直接把 $x=0$ 代入非零点求导所得的商式。
\end{exercise}
\begin{exercise}{259}{幂指函数的微分}
求 $y=(1+e^x)^x$ 的微分，逐项标明对底数与指数求导的贡献；再求 $\dd y|_{x=0}$，检查是否与原函数在 $0$ 处的线性近似一致。
\end{exercise}
\begin{exercise}{260}{由微分确定函数}
在各自定义区间内，求导数分别为 $x^2$、$\cos3x$、$1/(x\ln x)$、$e^x/(1+e^{2x})$ 的全部函数。写明任意常数与定义区间的关系，并直接求导核对四个结果。
\end{exercise}
\begin{exercise}{261}{夹逼条件与可微性}
设 $f$ 在 $0$ 的某邻域内有定义且 $|f(x)|\le x^2$。证明 $f$ 在 $0$ 处连续、可导并求 $f'(0)$；构造满足此条件而导函数在 $0$ 附近无界的函数，说明一点可导不能保证邻域内导数有界。
\end{exercise}
\begin{exercise}{262}{定点导数的结构化计算}
分别求：$f(x)=3x+\frac{x^2-1}{x^2+1}\arctan\frac{2x}{1+x}$ 的 $f'(0)$；$p(x)=x(x-1)\cdots(x-100)$ 的 $p'(100)$；$y=\frac{(x^2+1)\sqrt{3x+1}}{(x+1)^2e^{2x}}$ 的 $y'(1)$。每式选取合适方法，避免完全展开乘积。
\end{exercise}
\begin{exercise}{263}{等价差分给出的点值和导数}
设 $f$ 在 $x_0$ 处可导，且 $f(x_0+3h)-3f(x_0-h)\sim h$（$h\to0$）。求 $f(x_0)$、$f'(x_0)$，并用 $f(x_0+t)=f(x_0)+f'(x_0)t+o(t)$ 验证参数条件充分。
\end{exercise}
\begin{exercise}{264}{可导信息与幂指极限}
设 $f$ 在 $0$ 附近连续可导，$f(0)=1$、$f'(0)=4$。求 $\lim_{x\to0}[f(\sin x)]^{1/(2x)}$，证明底数最终为正；将指数改为 $1/(cx)$（$c\ne0$）时写出一般结果。
\end{exercise}
\begin{exercise}{265}{切线截距构造数列}
设曲线 $y=f(x)=1/(1+x^2)$ 在 $(n,f(n))$ 处的切线与 $x$ 轴相交于 $(\xi_n,0)$，其中 $n\ge1$。求 $\xi_n$，再求 $\lim_{n\to\infty}f(\xi_n)$ 及 $\lim_{n\to\infty}n^2f(\xi_n)$。
\end{exercise}
\begin{exercise}{266}{复合差商与法线}
设 $f$ 连续可导，且 $\lim_{x\to1}\frac{f(3x-1)+3}{x-1}=3$。求 $f(2)$、$f'(2)$，再求曲线在 $(2,f(2))$ 处的法线方程，说明复合函数内部系数的作用。
\end{exercise}
\begin{exercise}{267}{幂根极限与不可导点}
设 $f(x)=\lim_{n\to\infty}\frac{x^2}{(1+x^{2n})^{1/n}}$。求其分段表达式，讨论整个实轴上的连续性与可导性；对每个不可导点计算左右导数。
\end{exercise}
\begin{exercise}{268}{由连续延拓反求点值}
设 $f$ 连续可导且 $f(1)=0$、$f'(1)=1$。令 $g(x)=\begin{cases}\frac{f(e^{2x})+\ln(1-4x)}x,&x\ne0,\\b,&x=0,\end{cases}$ 在 $0$ 的充分小邻域内定义。求使 $g$ 在 $0$ 连续的 $b$，并说明现有条件是否足以保证 $g'(0)$ 存在。
\end{exercise}
\begin{exercise}{269}{极限生成分段函数的可导条件}
设 $f(x)=\lim_{n\to\infty}\frac{x^2e^{n(x-1)}+3ax+b}{e^{n(x-1)}+1}$。求分段表达式；确定使其在 $1$ 处可导的全部 $a,b$；在这些参数下求整个实轴上的 $f'$，并判断其连续性。
\end{exercise}
\begin{exercise}{270}{互补内函数的链式法则}
设 $f$ 可导，$y=f(\sin^2x)+f(\cos^2x)$。求 $y'$、$y'(0)$，再证明该函数以 $\pi/2$ 为周期；说明这是否保证最小正周期一定等于 $\pi/2$。
\end{exercise}
\begin{exercise}{271}{抽象函数的两层复合}
设 $f:\mathbb R\to\mathbb R$ 可导，$y=f(f(1+3x^2))$。求 $y'$ 与 $y'(0)$；若 $f$ 二阶可导，再求 $y''(0)$。
\end{exercise}
\begin{exercise}{272}{公共切线与复合平方}
曲线 $y=f(x)$ 与 $y=(3+\ln x)/x$ 在横坐标 $1$ 处的切线相同。设 $F(x)=[f(e^x)]^2$，求 $F'(0)$；写明公共切线同时给出的点值信息和导数信息。
\end{exercise}
\begin{exercise}{273}{二阶可导与极限反推}
设 $f$ 在 $0$ 处二阶可导，且 $\lim_{x\to0}\frac{e^{f(x)}-1}{1-\cos x}=2$。依次求 $f(0)$、$f'(0)$、$f''(0)$；说明为何不能直接假定 $f(x)$ 从二阶项开始。
\end{exercise}
\begin{exercise}{274}{复合二阶导数的正则性条件}
设 $f$ 在 $0$ 附近连续可导，$y=f(x^2)$。用导数定义求 $y''(0)$，结果用 $f'(0)$ 表示；若仅假设 $f$ 在邻域内可导，判断该结论是否必成立，并构造反例或给出证明。
\end{exercise}
\begin{exercise}{275}{导数方程与高阶通式}
设 $f$ 任意阶可导且 $f'=e^f$。求 $f^{(n)}$（$n\ge1$）关于 $f$ 的表达式，并用数学归纳法证明；若 $f(0)=0$，再求 $f^{(n)}(0)$。
\end{exercise}
\begin{exercise}{276}{高阶零点与定点导数}
设 $n$ 为正整数，$f(x)=(x+1)^ne^{-x^2}$。用莱布尼茨公式求 $f^{(n)}(-1)$；再求 $f^{(n+1)}(-1)$，说明绝大多数项消失的原因。
\end{exercise}
\begin{exercise}{277}{反函数的二阶导数}
设 $x=\varphi(y)$ 是 $y=x+\ln x$（$x>0$）的反函数。证明反函数在 $\mathbb R$ 上存在，求 $\dd^2x/\dd y^2$，并求其在 $y=1$ 处的值。
\end{exercise}
\begin{exercise}{278}{函数及导函数的逐层正则性}
设 $f(x)=\begin{cases}x^3\sin(1/x)+\sin x+1,&x\ne0,\\1,&x=0.\end{cases}$ 判断 $f$ 在 $0$ 处的连续性和可导性；求 $f'$，讨论其在 $0$ 处的连续性与可导性；求 $f''$ 的存在范围及表达式。
\end{exercise}
\begin{exercise}{279}{隐函数定点导数与无穷远极限}
方程 $\ln(x^2+y)=x^2y+\sin x$ 在 $(0,1)$ 附近确定分支 $y=f(x)$。求 $f'(0)$，再求 $\lim_{|x|\to\infty}x\left(1-f(2/x)\right)$；说明无穷远表达式只使用了 $f$ 在哪个点附近的信息。
\end{exercise}
\begin{exercise}{280}{隐函数与外层抽象函数的复合}
设 $f$ 二阶可导且 $f'(1)=1$，$xy-e^y+e^x=0$ 确定过原点的分支 $y=y(x)$。令 $z=f(e^{y(x)}-\sin x)$，求 $z''(0)$；计算中说明为何不需要知道 $f''(1)$ 的具体值。
\end{exercise}
\begin{exercise}{281}{同一横坐标对应的参数分支}
设 $x=2t^2+3t+1$、$y=e^t\sin t+2$。求 $x=0$ 对应的全部参数及各分支的 $\dd y/\dd x$；说明题目只给横坐标时，导数是否必为唯一数值。
\end{exercise}
\begin{exercise}{282}{参数函数的三阶导数}
设 $x=t-\ln\sqrt{1+e^{2t}}$、$y=\arctan e^t$。求 $\dd^3y/\dd x^3$，结果可用 $t$ 表示；证明 $\dd x/\dd t$ 恒非零，写出 $x$ 的值域。
\end{exercise}
\begin{exercise}{283}{极坐标曲线的切线与法线}
求阿基米德螺线 $\rho=\theta$ 在 $(\rho,\theta)=(\pi/2,\pi/2)$ 处的切线与法线方程，以直角坐标表示；先写出参数式并检查该点切线是否竖直。
\end{exercise}
\begin{exercise}{284}{线性主部反求复合导数}
设 $f$ 可微，$y=f(e^{2x})$ 在 $x=0$ 处，当 $\Delta x=-0.1$ 时，增量的线性主部也为 $-0.1$。求 $f'(1)$；说明这里给出线性主部是否意味着真实增量恰好等于 $-0.1$。
\end{exercise}
\begin{exercise}{285}{复合函数的定点微分}
设 $f'(x)=\arctan(1/x)$（$x\ne0$），$y=f\left(\frac{2x-1}{x+1}\right)$。求 $\dd y|_{x=2}$；列出复合求导公式需要排除的 $x$，说明分别来自哪一层函数。
\end{exercise}
\begin{exercise}{286}{参数函数的微分}
设 $f$ 可导且 $f'(0)\ne0$，$x=f(\ln t)$、$y=f(t^2-1)$（$t>0$）。求 $t=1$ 时的 $\dd y$，分别用 $\dd x$ 与 $\dd t$ 表示，并验证两种写法一致。
\end{exercise}
\begin{exercise}{287}{线性近似与误差控制}
用微分估算 $\sqrt{4.04}$，并利用函数在 $[4,4.04]$ 上的二阶导数给出误差上界；判断近似值偏大还是偏小。说明只有 $o(\Delta x)$ 信息时，能否直接得到给定步长下的数值误差上界。
\end{exercise}
\begin{exercise}{288}{隐函数微分与局部近似}
设方程 $ye^y=x$ 在 $(x,y)=(e,1)$ 附近确定函数 $y(x)$。求该点的 $\dd y$ 与 $y''$，写出 $x=e+h$ 时到二阶的近似式；说明为什么近似中出现的 $h^2$ 项不属于一阶微分。
\end{exercise}
